\documentclass[t]{beamer}
\usepackage{susu}
\usepackage{subcaption}
\DeclareCaptionSubType*{figure}
\renewcommand\thesubfigure{\asbuk{subfigure})}
\DeclareCaptionLabelFormat{russian}{{#1~#2}}
\captionsetup[subfigure]{labelformat=russian, font=small}
\begin{document}
	
	\firstpage
	\justifying
\section{Введение}

	\begin{frame}
		\frametitle{Цели и задачи} 
		\textbf{Цель:} разработка алгоритма сегментации факела выбросов с использованием тепло-видео систем.
		
		\textbf{Задачи:}
		\begin{enumerate}
			\justifying
			\item исследование существующих методов для решения задачи контроля выбросов;
			\item исследование современных способов применения оптических и тепловых снимков;
			\item разработка алгоритма для сегментации факела выбросов с использованием тепловых и оптических снимков.
		\end{enumerate}
	\end{frame}

\section[Современные методы]{Современные методы контроля выбросов}
	\begin{frame}
		\frametitle{\insertsection} 
		\begin{figure}
			\centering
			\includegraphics[width = \textwidth]{image/scheme1_upd}	
			\caption{Классификация методов\\ контроля выбросов}
			\label{fig:scheme1}
		\end{figure}
	\end{frame}

	\begin{frame}
		\frametitle{Преимущества тепловизоров}
		\begin{enumerate}
			\justifying
			\item большая дешевизна;
			\item относительно высокая точность.
		\end{enumerate}
		\begin{figure}[ht!]
			\begin{subfigure}{.45\textwidth}
				\centering
				\includegraphics[width = \textwidth]{image/optic_invis}
				\caption{}
			\end{subfigure}
			\begin{subfigure}{.45\textwidth}
				\centering
				\includegraphics[width = \textwidth]{image/tep_vis}
				\caption{}
			\end{subfigure}
			\centering
			\caption{Преимущество тепловых снимков перед\\ оптическими: a) оптическое изображение;\\ б) тепловизионное изображение}
			\label{fig:Examples}
		\end{figure}
	\end{frame}

\section{Задача}

	\begin{frame}
		\frametitle{Задача сегментации факела выбросов} 
		\vspace*{-0.3cm}
		Необходимо восстановить целевую функцию
		\begin{equation}
			f: X \rightarrow Z,
			\label{eq:segment_func}
		\end{equation}
	где $X$ -- пространство пар из RGB изображений и матриц температур;
	$Z$ -- пространство масок соответствующей размерности, где каждый пиксель отражает вероятность принадлежности к факелу.
	\begin{figure}[ht!]
		\begin{subfigure}{.25\textwidth}
			\centering
			\includegraphics[width = \textwidth]{image/opt_examp}
			\caption{}
		\end{subfigure}
		\begin{subfigure}{.25\textwidth}
			\centering
			\includegraphics[width = \textwidth]{image/tep_examp}
			\caption{}
		\end{subfigure}
		\begin{subfigure}{.25\textwidth}
			\centering
			\includegraphics[width = \textwidth]{image/mask}
			\caption{}
		\end{subfigure}
		\centering
		\vspace{-0.4cm}
		\caption{Примеры полученных изображений:\\ a) RGB изображение; б) матрица температур;\\ в) полученная маска}
		\label{fig:Examples}
	\end{figure}
	\end{frame}

	\begin{frame}
		\frametitle{Подзадачи поставленной задачи} 
		\vspace{1.8cm}
		\begin{figure}
			\centering
			\includegraphics[width = \textwidth]{image/scheme2_upd}
			\caption{Основные этапы поставленной задачи}
			\label{fig:scheme2_upd}
		\end{figure}
	\end{frame}

\section{Подготовка данных}
%	\begin{frame}
%		\frametitle{Работа с тепловизором}
%		\vspace{1cm}
%		\begin{figure}[ht!]
%			\begin{subfigure}{0.5\textwidth}
%				\centering
%				\includegraphics[width = \textwidth]{image/chapter_2/opt_example}
%				\caption{}
%			\end{subfigure}
%			\begin{subfigure}{0.3\textwidth}
%				\centering
%				\includegraphics[width = \textwidth]{image/chapter_2/tep_example}
%				\caption{}
%			\end{subfigure}
%			\centering
%			\caption{Примеры полученных изображений: \\a) оптический снимок; б) тепловой снимок}
%			\label{fig:Examples}
%		\end{figure}q
%	\end{frame}

%	\begin{frame}
%		\frametitle{Работа с тепловизором}
%		\vspace*{-0.45cm}
%		\begin{figure}[h!]
%			\centering
%			\includegraphics[width = \textwidth]{image/chapter_2/loaddata}	
%			\caption{Алгоритм сохранения снимков}
%			\label{fig:loaddata}
%		\end{figure}
%	\end{frame}

	\begin{frame}
		\frametitle{Преобразование цветовой карты}
		\vspace*{0.45cm}
		\begin{figure}[ht!]
			\begin{subfigure}{.45\textwidth}
				\centering
				\includegraphics[width = 4cm]{image/chapter_2/tep_example}
				\caption{}
			\end{subfigure}
			\begin{subfigure}{.45\textwidth}
				\centering
				\includegraphics[width = 4cm]{image/chapter_2/gray_tep_example}
				\caption{}
			\end{subfigure}
			\centering
			\caption{Примеры преобразования цветовой карты:\\ a) до; б) после}
			\label{fig:ResKNN}
		\end{figure}
	\end{frame}

	\begin{frame}
		\frametitle{Модель классификации <<k ближайших соседей>>}
		
		\vspace{0.85cm}
		\begin{figure}[h!]
			\centering
			\includegraphics[width = .8\textwidth]{image/cool_scheme}	
			\caption{Перевод цвета из RGB в оттенки\\ серого}
			\label{fig:kdtreeexample}
		\end{figure}
	\end{frame}
	
	\begin{frame}
		\frametitle{Алгоритм преобразования цветов}
		\vspace*{-0.45cm}
		\begin{figure}[ht!]
			\begin{subfigure}{.62\textwidth}
				\centering
				\includegraphics[width = \textwidth]{image/chapter_2/colorclassification}
				\vspace*{-0.6cm}
				\caption{}
			\end{subfigure}
			\begin{subfigure}{.2\textwidth}
				\centering
				\includegraphics[width = \textwidth]{image/chapter_2/colorclassification2}
				\vspace*{-0.6cm}
				\caption{}
			\end{subfigure}
			\centering
			\vspace*{-0.3cm}
			\caption{Алгоритмы преобразования цветов:\\ а) подготовка классификатора цветов; б) преобразование}
			\label{fig:Examples}
		\end{figure}
		
	\end{frame}

	\begin{frame}
		\frametitle{Метрика точности преобразования цветов}
%		\vspace*{0.2cm}
		\begin{equation}
			Acc = 1 - \frac{\sum\limits_{i=1}^h \sum\limits_{j=1}^w |P^{true}_{ij} - P^{conv}_{ij}|}{255wh},
			\label{eq:flanaccuracy}	
			\vspace*{0.2cm}
		\end{equation}
		где $w$ -- высота кадра;\\ \hfill \break
		\hspace*{0.65cm}$h$ -- ширина кадра;\\ \hfill \break
		\hspace*{0.65cm}$P^{true}$ -- изображение в оттенках серого;\\ \hfill \break
		\hspace*{0.65cm}$P^{conv}$ -- изображение, с наложенной цветовой картой, сжатое с помощью JPEG и обработанное алгоритмом преобразования. 
	\end{frame}

	\begin{frame}
		\frametitle{Метрика точности преобразования цветов}
		\vspace*{-0.35cm}
		\begin{figure}[h!]
			\centering
			\includegraphics[width = 0.8\textwidth]{image/accuracy_plot}	
			\caption{Зависимость точности от степени сжатия}
			\label{fig:fullprepare}
		\end{figure}
	
	\end{frame}

	\begin{frame}
		\frametitle{Алгоритм подготовки данных}
		\vspace*{-0.2cm}
					\begin{figure}[ht!]
						\hspace*{-1.325cm}
							\begin{subfigure}{0.4\textwidth}
									\centering
									\includegraphics[width = \textwidth]{image/examp111}
							\end{subfigure}
							\begin{subfigure}{0.3\textwidth}
								\centering
								\includegraphics[width = \textwidth]{image/1}
							\end{subfigure}
							
							\begin{subfigure}{0.3\textwidth}
									\centering
									\includegraphics[width = \textwidth]{image/examp222}
									\caption{}
							\end{subfigure}
							\begin{subfigure}{0.3\textwidth}
								\centering
								\includegraphics[width = \textwidth]{image/1tep}
								\caption{}
							\end{subfigure}
							\centering
							\vspace*{-0.25cm}
							\caption{Примеры подготовки данных: \\a) снимки, полученые с тепловизора; б) снимки,\\ подготовленые к сегментации}
							\label{fig:Examples}
						\end{figure}

		
	\end{frame}

\section{Сегментация}

	\begin{frame}
		\frametitle{Алгоритм детекции трубы}
		\vspace*{0.6cm}
		\begin{figure}[h!]
			\centering
			\includegraphics[width = \textwidth]{image/match1}
			\caption{Визуализация работы алгоритма детекции трубы}
			\label{fig:fullprepare}
		\end{figure}
	\end{frame}	

	\begin{frame}
		\frametitle{Алгоритм детекции трубы}
		\vspace*{-0.2cm}
		\begin{figure}[ht!]
			\begin{subfigure}{.51\textwidth}
				\centering
				\includegraphics[width = \textwidth]{image/trubadetection}
				\vspace*{-0.6cm}
				\caption{}
			\end{subfigure}
			\begin{subfigure}{.46\textwidth}
				\centering
				\includegraphics[width = \textwidth]{image/RanSaC}
				\vspace*{-0.6cm}
				\caption{}
			\end{subfigure}
			\centering
			\vspace*{-0.3cm}
			\caption{Алгоритмы детекции трубы:\\ а) поиск ключевых точек; б) восстановление прямоугольника}
			\label{fig:Examples}
		\end{figure}
		
	\end{frame}
	

\begin{frame}
	\frametitle{Алгоритм сегментации методом водораздела}
	\vspace*{0.7cm}
	\begin{figure}[ht!]
		\begin{subfigure}{0.32\textwidth}
			\centering
			\includegraphics[width = \textwidth]{image/WaterShed1}
			\caption{}
		\end{subfigure}
		\begin{subfigure}{0.32\textwidth}
			\centering
			\includegraphics[width = \textwidth]{image/WaterShed2}
			\caption{}
		\end{subfigure}
		\begin{subfigure}{0.32\textwidth}
			\centering
			\includegraphics[width = \textwidth]{image/WaterShed3}
			\caption{}
		\end{subfigure}
		\caption{Пример построение водораздела: а) функция от координат пикселя; б) процесс заполнения <<водой>>;\\ в) построенный водораздел}
		\label{fig:WaterShed}
	\end{figure}
\end{frame}

	\begin{frame}
		\frametitle{Алгоритм сегментации методом водораздела}
		\vspace*{0.7cm}
		\begin{figure}[h!]
			\centering
			\includegraphics[width = 0.6\textwidth]{image/waterShedSegmentation}
			\caption{Алгоритм сегментации\\ методом водораздела}
			\label{fig:fullprepare}	
		\end{figure}
	\end{frame}

\section{Результаты}

	\begin{frame}
		\frametitle{Результаты работы алгоритма сегментации}
		\vspace*{-0.35cm}
		\begin{figure}[h!]
			\begin{subfigure}{.30\textwidth}
				\centering
				\includegraphics[width = \textwidth]{image/examples/img/206}
				%			\caption{}
			\end{subfigure}
			\begin{subfigure}{.30\textwidth}
				\centering
				\includegraphics[width = \textwidth]{image/examples/tep/206}
				%			\caption{}
			\end{subfigure}
			\begin{subfigure}{.30\textwidth}
				\centering
				\includegraphics[width = \textwidth]{image/examples/mask/206}
				%			\caption{}
			\end{subfigure}
			\hspace*{0.00cm}
			\begin{subfigure}{.30\textwidth}
				\centering
				\includegraphics[width = \textwidth]{image/examples/img/240}
				\caption{}
			\end{subfigure}
%			\hspace{0.1cm}
			\begin{subfigure}{.30\textwidth}
				\centering
				\includegraphics[width = \textwidth]{image/examples/tep/240}
				\caption{}
			\end{subfigure}
%			\hspace{0.1cm}
			\begin{subfigure}{.30\textwidth}
				\centering
				\includegraphics[width = \textwidth]{image/examples/mask/240}
				\caption{}
			\end{subfigure}
			\vspace{-0.3cm}
			\caption{Пример работы: а) оригинальные оптические изображения; б) оригинальные тепловые изображения;\\ в) получившиеся маски}
			\label{fig:test_segm}
		\end{figure}
	\end{frame}

	\begin{frame}
		\frametitle{Разница масок}
		\vspace*{-0.35cm}
		\begin{figure}[h!]
			\begin{subfigure}{.30\textwidth}
				\centering
				\includegraphics[width = \textwidth]{image/examples/mask/206}
				%			\caption{}
			\end{subfigure}
			\begin{subfigure}{.30\textwidth}
				\centering
				\includegraphics[width = \textwidth]{image/examples/mask_razmet/206}
				%			\caption{}
			\end{subfigure}
			\begin{subfigure}{.30\textwidth}
				\centering
				\includegraphics[width = \textwidth]{image/examples/mask_dif/206}
				%			\caption{}
			\end{subfigure}
			\hspace*{0.00cm}
			\begin{subfigure}{.30\textwidth}
				\centering
				\includegraphics[width = \textwidth]{image/examples/mask/240}
				\caption{}
			\end{subfigure}
			%			\hspace{0.1cm}
			\begin{subfigure}{.30\textwidth}
				\centering
				\includegraphics[width = \textwidth]{image/examples/mask_razmet/240}
				\caption{}
			\end{subfigure}
			%			\hspace{0.1cm}
			\begin{subfigure}{.30\textwidth}
				\centering
				\includegraphics[width = \textwidth]{image/examples/mask_dif/240}
				\caption{}
			\end{subfigure}
			\vspace{-0.3cm}
			\caption{Сравнение масок: а) результат работы алгоритма;\\ б) размеченная маска; в) разница масок}
			\label{fig:test_segm}
		\end{figure}
	\end{frame}

	\begin{frame}
		\frametitle{Метрика точности сегментации}
		\vspace{0.5cm}
		\begin{equation}
			\textup{DICE} = \frac{2\textup{TP}}{2\textup{TP} + \textup{FP} + \textup{FN}},
			\label{eq:DiceLoss}		\vspace{0.3cm}
		\end{equation}

		где TP -- правильное предсказание сегмента;\\ \hfill \break
		\hspace*{0.65cm}FP -- ложное предсказание сегмента;\\ \hfill \break
		\hspace*{0.65cm}FN -- неправильное предсказание сегмента.\\ \hfill \break
		
		Итоговая точность -- 86,2\%
	\end{frame}

	\begin{frame}
		\frametitle{Заключение}
		В ходе работы разработан алгоритм сегментации факела выбросов, использующий данные тепло-видео систем наблюдения. Помимо этого:
		\begin{enumerate}
			\justifying
			\item проведено исследование существующих методов для решения задачи контроля выбросов;
			\item проведено исследование современных способов применения оптических и тепловых снимков;
			\item разработан алгоритм для сегментации факела выбросов с использованием тепловых и оптических снимков.
		\end{enumerate}
		
		
		\textit{Данная работа была выполнена в рамках проекта <<Экомонитор>>, Института естественных и точных наук.}
	\end{frame}

	\begin{frame}
		\vspace{2.9cm}
		\begin{center}{\Huge{Спасибо за внимание!}}\end{center}
	\end{frame}

\appendix
%
%\addtocounter{page}{-1}


	\begin{frame}[noframenumbering]
			\vspace*{0.7cm}
				\vspace*{0.7cm}
		\begin{figure}[h!]
			\centering
			\includegraphics[width = 0.6\textwidth]{image/chapter_2/jet}
			\caption{Цветовая карта <<JET>>}
			\label{fig:jet}	
		\end{figure}
	\end{frame}

\appendix
%\addtocounter{page}{-1}

\begin{frame}[noframenumbering]
		\vspace*{0.7cm}
	\begin{figure}[h!]
		\centering
		\includegraphics[width = 0.6\textwidth]{image/accuracy_itog_plot}
		\caption{График зависимости точности от номера кадра}
		\label{fig:accuracy_itog_plot}	
	\end{figure}
\end{frame}

\appendix
%\addtocounter{page}{-1}

\begin{frame}[noframenumbering]
		\vspace*{0.7cm}	\vspace*{0.7cm}
	\begin{figure}[h!]
		\begin{subfigure}{.28\textwidth}
			\centering
			\includegraphics[width = 1\textwidth]{image/chapter_2/SKIZ_Watershed_build1}
			\caption{}
		\end{subfigure}
		\begin{subfigure}{.20\textwidth}
			\centering
			\includegraphics[width = \textwidth]{image/chapter_2/SKIZ_Watershed_build2}
			\caption{}
		\end{subfigure}
		\centering
		\caption{Пример построения водораздела:\\ a) секции уровня $Z_i$ и $Z_{i+1}$; б) построенный водораздел}
		\label{fig:SKIZ_Watershed_build}
	\end{figure}
\end{frame}

\appendix

\begin{frame}[noframenumbering]
		\begin{figure}[h!]
		\begin{subfigure}{.30\textwidth}
			\centering
			\includegraphics[width = \textwidth]{image/chapter_3/examples/img/159}
			%			\caption{}
		\end{subfigure}
		\begin{subfigure}{.30\textwidth}
			\centering
			\includegraphics[width = \textwidth]{image/chapter_3/examples/tep/159}
			%			\caption{}
		\end{subfigure}
		\begin{subfigure}{.30\textwidth}
			\centering
			\includegraphics[width = \textwidth]{image/chapter_3/examples/mask/159}
			%			\caption{}
		\end{subfigure}
		\begin{subfigure}{.30\textwidth}
			\centering
			\includegraphics[width = \textwidth]{image/chapter_3/examples/img/214}
			\caption{}
		\end{subfigure}
%		\hspace{0.1cm}
		\begin{subfigure}{.30\textwidth}
			\centering
			\includegraphics[width = \textwidth]{image/chapter_3/examples/tep/214}
			\caption{}
		\end{subfigure}
%		\hspace{0.1cm}
		\begin{subfigure}{.30\textwidth}
			\centering
			\includegraphics[width = \textwidth]{image/chapter_3/examples/mask/214}
			\caption{}
		\end{subfigure}
%		\vspace*{-0.5cm}
		\caption{Пример работы: а) оригинальные оптические изображения; б) оригинальные тепловые изображения; в) получившиеся маски}
		\label{fig:test_segm}
	\end{figure}
\end{frame}
\end{document}
